\section{Arbitrary-spin onset on finite unicyclic networks}
\label{sec:unicyclic}

We now determine the first coefficient allowed by gauge invariance and physical
time reversal on a finite unicyclic graph.  The result applies to arbitrary
irreducible site spins.  It also allows arbitrary longitudinal exchanges,
longitudinal fields, and trees grafted onto the unique cycle.  All Hamiltonian
parameters are held fixed as $\beta\to0$.

Let $G=(V,E)$ be finite, simple, connected, and unicyclic.  Denote its unique
cycle by $\gamma$ and its length by $n\geq3$.  Choose three distinct vertices
$A,B,C\in V(\gamma)$.  Orient $\gamma$ so that they occur in the cyclic order

\begin{equation}
 \gamma=(A,\ldots,B,\ldots,C,\ldots,A).
 \label{eq:marked-cycle-order}
\end{equation}

For this orientation, set

\begin{equation}
 W_\gamma=\prod_{(x,y)\in\gamma}t_{xy}.
 \label{eq:unicyclic-W}
\end{equation}

The notation $AB$ below means the subsystem $\{A,B\}$.  It does not assert
that $A$ and $B$ are adjacent.  The same convention applies to $BC$.

\begin{theorem}[Arbitrary-spin unicyclic onset]
\label{thm:unicyclic-onset}
For the axial Hamiltonian \eqref{eq:axial-H}, let every $\mathcal H_x$ carry
the irreducible spin-$s_x$ representation, with $\kappa_x$ defined in
Eq.~\eqref{eq:kappa-definition}.  Assume $t_{xy}\in\mathbb C$ with
$t_{yx}=\overline{t_{xy}}$, $J^z_{xy}\in\mathbb R$, and $h_x\in\mathbb R$.
Then every Taylor coefficient of
$\Mmod_\beta(AB,BC)$ through order $n$ vanishes, and

\begin{equation}
 \boxed{
 \Mmod_\beta(AB,BC)
 =2(-1)^n\beta^{n+1}h_B
 \left(\prod_{x\in V(\gamma)}\kappa_x\right)
 \operatorname{Im}W_\gamma
 +O(\beta^{n+2}) .}
 \label{eq:unicyclic-onset}
\end{equation}

No coupling is assumed to be nonzero.  Equation~\eqref{eq:unicyclic-onset}
is therefore a statement about the leading admissible coefficient.  Its
displayed coefficient can vanish for a particular Hamiltonian.
\end{theorem}

The proof occupies the rest of this section.  We first isolate the only
possible monomial at order $n+1$.  We then evaluate its coefficient by
reducing the cycle to three primitive arc cumulants.  This second step is
needed because tracing out an arc generates multispin operators in addition
to its endpoint hopping.

\subsection{Gauge-flow grading and pruning}

Write the analytic expansion near the maximally mixed state as

\begin{equation}
 \Mmod_\beta(AB,BC)=\sum_{q\geq0}\beta^q m_q.
 \label{eq:unicyclic-series}
\end{equation}

Finite dimensionality makes this series convergent in a neighborhood of
$\beta=0$.  The identity $\rho_\beta(H)=\rho_1(\beta H)$ shows that $m_q$ is
homogeneous of total degree $q$ in the Hamiltonian parameters.

For coefficient extraction, first write $t_e=x_e+\ii y_e$.  Each $m_q$ is a
real polynomial in $x,y,J^z,h$.  Complexify this polynomial and make the
invertible change of variables

\begin{equation}
 z_e=x_e+\ii y_e,
 \qquad
 w_e=x_e-\ii y_e.
 \label{eq:formal-complexification}
\end{equation}

The variables $z_e$ and $w_e$ are now algebraically independent.  The
physical coefficient is recovered on $w_e=\overline{z_e}$.  Gauge invariance
under the vertex torus is a polynomial identity, so distinct torus characters
vanish coefficientwise.  Equations~\eqref{eq:M-K-odd} and
\eqref{eq:M-Theta-odd} give the polynomial identities

\begin{align}
 m_q(w,z,J^z,h)&=-m_q(z,w,J^z,h),
 \label{eq:formal-K-oddness}\\
 m_q(z,w,J^z,-h)&=-m_q(z,w,J^z,h).
 \label{eq:formal-Theta-oddness}
\end{align}

This formalization also makes every edge multidegree unambiguous.

Fix an orientation of each edge.  For a monomial, let $a_e$ and $b_e$ be the
exponents of $z_e$ and $w_e$, and define the integral edge flow

\begin{equation}
 d_e=a_e-b_e.
 \label{eq:edge-flow}
\end{equation}

Local rotations about the $z$ axis leave the modular commutator invariant.
Character independence under this torus action means that the coefficient of
a monomial can be nonzero only if

\begin{equation}
 \operatorname{div}d=0
 \label{eq:flow-divergence}
\end{equation}

at every vertex.  On a unicyclic graph, any integral flow satisfying
Eq.~\eqref{eq:flow-divergence} vanishes on every bridge.  Its restriction to
$\gamma$ is $k$ times the consistently oriented unit flow for some
$k\in\mathbb Z$.  This follows by removing the cycle edges, starting at a
leaf of each remaining tree, and imposing zero divergence inward.  Zero
divergence on the cycle then makes all oriented cycle values equal.

Let $D_{\mathrm{hop}}=\sum_e(a_e+b_e)$ be the hopping degree.
Equation~\eqref{eq:formal-K-oddness} removes the sector $d=0$.  Indeed, every
monomial in that sector is fixed by $z\leftrightarrow w$ after identical
factors are collected.  Every surviving monomial thus has $k\neq0$, and

\begin{equation}
 D_{\mathrm{hop}}
 \geq\sum_e|d_e|
 =n|k|
 \geq n.
 \label{eq:hopping-lower-bound}
\end{equation}

Equation~\eqref{eq:formal-Theta-oddness} says that the total field degree of
every surviving monomial is odd and in particular is at least one.  In the
physical Hamiltonian, this identity follows because time reversal leaves $t$
and $J^z$ fixed in the parameterization used here, sends $h$ to $-h$, and
reverses the modular commutator.  Equations~\eqref{eq:unicyclic-series} and
\eqref{eq:hopping-lower-bound} now imply

\begin{equation}
 m_q=0\qquad(0\leq q\leq n).
 \label{eq:unicyclic-low-orders-zero}
\end{equation}

Equality in the degree bound at order $n+1$ is rigid.  The flow is one unit
around $\gamma$, one oriented hopping monomial occurs on every cycle edge,
and exactly one longitudinal field occurs.  There is no degree
left for a longitudinal exchange, a bridge hopping, or a neutral return.
After combining the two conjugate orientations, one obtains

\begin{equation}
 m_{n+1}=\operatorname{Im}W_\gamma\sum_{x\in V}c_xh_x
 \label{eq:unicyclic-leading-classification}
\end{equation}

with real constants $c_x$ that can a priori depend on the graph and its site
spins.

Equation~\eqref{eq:unicyclic-leading-classification} also proves a useful
pruning statement.  Extract the coefficient of $h_xW_\gamma$ and set every
parameter absent from that monomial to zero.  Every grafted tree then becomes
a decoupled tensor factor.  If $x\notin V(\gamma)$, the remaining field acts
on a factor disjoint from $ABC$, so the reduced state on $ABC$ is independent
of $h_x$ and $c_x=0$.  If $x\in V(\gamma)$, every off-cycle normalized trace
equals one.  The coefficient is therefore exactly the coefficient of the
isolated marked cycle, and the remaining constants depend only on the cycle
spins.  We may work with $H_\gamma$ from now on.

\subsection{Three-site expansion through fourth order}
\label{sec:three-site-expansion}

The triangular seed below uses a general local expansion that we record here.
In this subsection, $H$ is a Hamiltonian on
$\mathcal H_A\otimes\mathcal H_B\otimes\mathcal H_C$.  A scalar can be
subtracted from $H$ without changing its Gibbs state, so we take
$\tau(H)=0$, where

\begin{equation}
 \tau(O)=\frac1d\operatorname{Tr}O,
 \qquad
 d=\dim(\mathcal H_A\otimes\mathcal H_B\otimes\mathcal H_C).
 \label{eq:normalized-trace}
\end{equation}

For $X\subseteq ABC$, let $X^c$ denote the complement inside $ABC$ and
introduce the trace-preserving conditional expectation

\begin{equation}
 \E_X(O)
 =\frac{1}{d_{X^c}}\operatorname{Tr}_{X^c}(O)\otimes I_{X^c}.
 \label{eq:conditional-expectation}
\end{equation}

Set

\begin{equation}
 r_\beta=d\rho_\beta
 =\frac{e^{-\beta H}}{\tau(e^{-\beta H})},
 \qquad
 L_X(\beta)=\log\E_X(r_\beta).
 \label{eq:r-and-L}
\end{equation}

Since
$\E_X(r_\beta)=d_X\rho_{\beta,X}\otimes I_{X^c}$,
$L_X$ differs from the extended operator
$\widehat{\log\rho_{\beta,X}}$ only by a scalar.  Hence

\begin{equation}
 \Mmod_\beta(AB,BC)
 =\ii\tau\left(r_\beta
 [L_{AB}(\beta),L_{BC}(\beta)]\right).
 \label{eq:M-normalized}
\end{equation}

The normalized Gibbs density has the expansion

\begin{equation}
 r_\beta
 =I-\beta H
 +\frac{\beta^2}{2}
 \left(H^2-\tau(H^2)I\right)
 +O(\beta^3).
 \label{eq:r-expansion}
\end{equation}

Define

\begin{equation}
 P=\E_{AB}(H),
 \qquad
 Q=\E_{BC}(H),
 \label{eq:PQ}
\end{equation}

and, modulo scalar multiples of the identity,

\begin{align}
 A_2&=\frac12\left(\E_{AB}(H^2)-P^2\right),
 \label{eq:A2}\\
 B_2&=\frac12\left(\E_{BC}(H^2)-Q^2\right).
 \label{eq:B2}
\end{align}

Expanding the logarithms gives

\begin{align}
 L_{AB}(\beta)&=-\beta P+\beta^2A_2+O(\beta^3),
 \label{eq:LAB-expansion}\\
 L_{BC}(\beta)&=-\beta Q+\beta^2B_2+O(\beta^3).
 \label{eq:LBC-expansion}
\end{align}

\begin{lemma}[Expansion through fourth order]
\label{lem:fourth-order-expansion}
The coefficient of $\beta^2$ in the modular commutator vanishes.  The next
two coefficients are

\begin{align}
 [\beta^3]\Mmod_\beta(AB,BC)
 &=-\ii\tau\bigl(H[P,Q]\bigr),
 \label{eq:m3-general}\\
 [\beta^4]\Mmod_\beta(AB,BC)
 &=\ii\tau\left[
 \frac12H^2[P,Q]
 +H\bigl([P,B_2]+[A_2,Q]\bigr)
 \right].
 \label{eq:m4-general}
\end{align}
\end{lemma}

\begin{proof}
Insert Eqs.~\eqref{eq:r-expansion} through
\eqref{eq:LBC-expansion} into Eq.~\eqref{eq:M-normalized}.  At order
$\beta^2$ one obtains $\ii\tau([P,Q])=0$.  At order $\beta^3$, every bare
commutator has zero trace, leaving Eq.~\eqref{eq:m3-general}.  At order
$\beta^4$, terms involving the third-order coefficients of $L_{AB}$ or
$L_{BC}$ again occur as bare commutators.  The remaining contributions are
those in Eq.~\eqref{eq:m4-general}.
\end{proof}

\subsection{Local spin moments and the triangular seed}
\label{sec:arbitrary-spin-triangle}

Let

\begin{equation}
 \tau_x(O)=\frac{1}{2s_x+1}\operatorname{Tr}_x O.
 \label{eq:local-normalized-trace}
\end{equation}

Rotational invariance of the normalized trace and the spin commutation
relations give

\begin{align}
 \tau_x(S_x^a)&=0,
 \label{eq:spin-first-moment}\\
 \tau_x(S_x^aS_x^b)&=\kappa_x\delta_{ab},
 \label{eq:spin-second-moment}\\
 \tau_x(S_x^aS_x^bS_x^c)
 &=\frac{\ii\kappa_x}{2}\varepsilon_{abc}.
 \label{eq:spin-third-moment}
\end{align}

For completeness, the second moment is an invariant rank-two tensor.  Its
trace over $a=b$ is $s_x(s_x+1)$, which gives
Eq.~\eqref{eq:spin-second-moment}.  The invariant rank-three tensor is
antisymmetric.  Taking the trace of
$[S_x^a,S_x^b]S_x^c=\ii\varepsilon_{abd}S_x^dS_x^c$ fixes its coefficient
and proves Eq.~\eqref{eq:spin-third-moment}.  In the ladder basis these
identities read

\begin{equation}
 \tau_x(S_x^+S_x^-)=\tau_x(S_x^-S_x^+)=2\kappa_x,
 \label{eq:spin-ladder-second-moment}
\end{equation}

while the six orderings of $S_x^+,S_x^-,S_x^z$ have normalized trace
$+\kappa_x$ or $-\kappa_x$, three of each.

\begin{proposition}[Arbitrary-spin triangular seed]
\label{prop:arbitrary-spin-triangle}
For the isolated three-site axial triangle oriented as
$A\to B\to C\to A$, with arbitrary irreducible spins and arbitrary axial
couplings,

\begin{equation}
 \Mmod_\beta(AB,BC)
 =-2\beta^4h_B\kappa_A\kappa_B\kappa_C
 \operatorname{Im}(t_{AB}t_{BC}t_{CA})+O(\beta^5).
 \label{eq:arbitrary-spin-triangle}
\end{equation}

The coefficient is independent of $h_A$, $h_C$, and all $J^z$ couplings.
\end{proposition}

\begin{proof}
Homogeneity, gauge invariance, conjugation oddness, and field oddness give

\begin{equation}
 [\beta^4]\Mmod_\beta(AB,BC)
 =(c_Ah_A+c_Bh_B+c_Ch_C)
 \operatorname{Im}(t_{AB}t_{BC}t_{CA}).
 \label{eq:triangle-arbitrary-classification}
\end{equation}

There is no degree available for $J^z$.  Evaluate the three constants at the
unit-amplitude parameter point

\begin{equation}
 t_{AB}=t_{BC}=1,
 \qquad
 t_{CA}=\ii,
 \qquad
 J^z=0.
 \label{eq:arbitrary-triangle-point}
\end{equation}

Use $P,Q,A_2,B_2$ from Eqs.~\eqref{eq:PQ} through \eqref{eq:B2} and write

\begin{align}
 T_0&=\frac{\ii}{2}\tau\bigl(H^2[P,Q]\bigr),
 \label{eq:triangle-T0}\\
 T_1&=\ii\tau\bigl(H[P,B_2]\bigr),
 \label{eq:triangle-T1}\\
 T_2&=\ii\tau\bigl(H[A_2,Q]\bigr).
 \label{eq:triangle-T2}
\end{align}

Lemma~\ref{lem:fourth-order-expansion} gives
$[\beta^4]\Mmod_\beta=T_0+T_1+T_2$.  Differentiation with respect to one
field at zero field gives the exact contraction table

\begin{equation}
\begin{array}{c|ccc|c}
 r&\partial_{h_r}T_0&\partial_{h_r}T_1&\partial_{h_r}T_2
 &\partial_{h_r}m_4\\
 \hline
 A&0&0&0&0\\
 B&2K_\triangle&-2K_\triangle&-2K_\triangle&-2K_\triangle\\
 C&0&0&0&0
\end{array}
\qquad
K_\triangle=\kappa_A\kappa_B\kappa_C.
 \label{eq:arbitrary-spin-contraction-table}
\end{equation}

We justify the lift from the explicit spin-$1/2$ contraction in
Appendix~\ref{app:pauli-certificate}.  In every nonzero square-free
contraction of $h_r t_{AB}t_{BC}t_{CA}$, each vertex distinct from $r$
carries one $S^+$ and one $S^-$.  The vertex $r$ carries one each of
$S^+$, $S^-$, and $S^z$.  If a conditional expectation splits the letters
at a vertex among two normalized traces, one factor contains a single
traceless generator and the contraction vanishes.  Every surviving
contraction therefore uses exactly one of
Eqs.~\eqref{eq:spin-ladder-second-moment} and
\eqref{eq:spin-third-moment} at each vertex.  Its ratio to the spin-$1/2$
value is $4\kappa_x$ at vertex $x$.  Multiplying every entry of the qubit
table by
$\prod_{x\in\{A,B,C\}}4\kappa_x$ gives
Eq.~\eqref{eq:arbitrary-spin-contraction-table}.  Summing its three columns
proves Eq.~\eqref{eq:arbitrary-spin-triangle}.
\end{proof}

This argument lifts the square-free contraction, not the Pauli
anticommutator identities.  Only the second and third normalized spin
moments are used.

\subsection{Primitive cumulants of the three marked arcs}

Split the oriented cycle into the three internally disjoint arcs

\begin{equation}
 P_1:A\longrightarrow B,
 \qquad
 P_2:B\longrightarrow C,
 \qquad
 P_3:C\longrightarrow A.
 \label{eq:three-marked-arcs}
\end{equation}

Let their lengths be $\ell_1,\ell_2,\ell_3\geq1$, so that
$\ell_1+\ell_2+\ell_3=n$.  Put $U=\{A,B,C\}$ and
$I=V(\gamma)\setminus U$.  With

\begin{equation}
 \tau_I(O)=\frac{1}{d_I}\operatorname{Tr}_I O,
 \qquad
 d_I=\prod_{x\in I}(2s_x+1),
 \label{eq:internal-normalized-trace}
\end{equation}

the normalized partial trace over $I$ defines

\begin{align}
 G(\beta)&=\tau_I(e^{-\beta H_\gamma}),
 \label{eq:arc-G}\\
 Q(\beta)&=-\log G(\beta).
 \label{eq:arc-Q}
\end{align}

The physical reduced state on $U$ is
$G/\operatorname{Tr}_UG$.  Define the real-analytic functional near $Q=0$

\begin{equation}
 \mathcal F(Q)
 =\Mmod_{e^{-Q}/\operatorname{Tr}_Ue^{-Q}}(AB,BC).
 \label{eq:triangle-functional}
\end{equation}

Then $\Mmod_\beta(AB,BC)=\mathcal F(Q(\beta))$.  The dimensionless
definition \eqref{eq:arc-Q} is essential.  The composed functional is
$\mathcal F(Q)$, not the same functional evaluated on
$-\beta^{-1}\log G$.

Consider an oriented arc

\begin{equation}
 P=(v_0=u,v_1,\ldots,v_\ell=v),
 \qquad
 T_P=\prod_{j=0}^{\ell-1}t_{v_jv_{j+1}},
 \label{eq:oriented-arc}
\end{equation}

and let $P^\circ=\{v_1,\ldots,v_{\ell-1}\}$.  The symbol $[P]$ denotes
extraction of the primitive multiaffine coefficient in which every oriented
hopping of $P$ occurs once and no other Hamiltonian parameter occurs.  The
monomial $T_P$ is retained by this notation.

\begin{lemma}[Primitive arc jet]
\label{lem:primitive-arc-jet}
For every marked arc $P$,

\begin{equation}
 [P]G
 =(-\beta)^\ell
 \left(\prod_{x\in P^\circ}\kappa_x\right)
 \frac{T_P}{2}S_u^+S_v^-.
 \label{eq:primitive-arc-G}
\end{equation}

Consequently,

\begin{equation}
 [P]Q=-[P]G.
 \label{eq:primitive-arc-Q}
\end{equation}
\end{lemma}

\begin{proof}
Only the term of degree $\ell$ in the exponential can contribute to
$[P]G$.  The bond normalization gives $2^{-\ell}$.  At every internal
vertex the two possible local orders are $S^-S^+$ and $S^+S^-$.  Both have
normalized trace $2\kappa_x$ by
Eq.~\eqref{eq:spin-ladder-second-moment}.  All $\ell!$ global orderings
therefore give the same partially traced operator, and they cancel the
factorial in the exponential.  The resulting scalar is
$2^{-\ell}\prod_{x\in P^\circ}2\kappa_x
=\frac12\prod_{x\in P^\circ}\kappa_x$, which proves
Eq.~\eqref{eq:primitive-arc-G}.

To extract $[P]\log G$, expand the logarithm around $G=I$.  Every nonlinear
term would split the edge set of $P$ into two or more nonempty proper
subsets.  Lemma~\ref{lem:proper-subpath}, proved next, makes each such factor
zero.  Thus $[P]\log G=[P]G$, and the minus sign in
Eq.~\eqref{eq:arc-Q} proves Eq.~\eqref{eq:primitive-arc-Q}.
\end{proof}

\begin{lemma}[Proper-subpath vanishing]
\label{lem:proper-subpath}
In the primitive edge sector of the three marked arcs, a coefficient block
means one factor arising in an ordered product from the exponential or
logarithmic expansion.  Such a coefficient of $G$ vanishes if its selected
edges contain a nonempty proper subset of at least one arc and do not contain
the rest of that arc in the same block.  The conclusion remains true when
the block contains edges from the other arcs and the unique longitudinal
field.
\end{lemma}

\begin{proof}
The selected edges of a proper arc subset have a connected component with an
endpoint in the interior of that arc.  At this endpoint the local operator
has axial charge $+1$ or $-1$.  Every other selected edge in the same block
lies either on the same proper subset or on an arc with a disjoint interior,
so it cannot cancel this local charge.  A longitudinal field has charge zero.
The normalized trace of a charged local operator vanishes under rotations
about the $z$ axis.  Hence the entire partial trace is zero.
\end{proof}

\begin{lemma}[Pairwise shuffle cancellation]
\label{lem:pairwise-shuffle}
Let $P_i$ and $P_j$ be two complete marked arcs.  In the primitive sector
with no further insertion,

\begin{equation}
 [P_iP_j]G
 =\frac12\bigl([P_i]G[P_j]G+[P_j]G[P_i]G\bigr).
 \label{eq:two-arc-shuffle}
\end{equation}

The same identity holds for a complete arc and one field on a retained
vertex.  It follows that

\begin{equation}
 [P_iP_j]Q=0,
 \qquad
 [P_i h_r]Q=0
 \quad(r\in U).
 \label{eq:pair-jets-zero}
\end{equation}

If the field is on an internal vertex of $P_i$, then

\begin{equation}
 [P_j h_r]Q=0
 \qquad(j=1,2,3)
 \label{eq:internal-field-pair-zero}
\end{equation}

as well.
\end{lemma}

\begin{proof}
The raw coefficient of two complete arcs is the sum over all shuffles of
their bond letters.  The internal traces are independent of the order by
Eq.~\eqref{eq:spin-ladder-second-moment}.  Exactly half of the shuffles put
the endpoint letter of $P_i$ before the endpoint letter of $P_j$, and half
put it after.  This proves Eq.~\eqref{eq:two-arc-shuffle}.  The argument for
an arc and a retained field is identical.

By Lemma~\ref{lem:proper-subpath}, the quadratic term is the only nonlinear
term in $\log G$ that can contribute to either two-ingredient coefficient.
It is exactly one half of the symmetrized product in
Eq.~\eqref{eq:two-arc-shuffle}.  Explicitly,

\begin{align}
 [P_iP_j]Q
 ={}&-[P_iP_j]G
 +\frac12\bigl([P_i]G[P_j]G+[P_j]G[P_i]G\bigr)
 =0.
 \label{eq:two-arc-Q-cancellation}
\end{align}

The same formula with one arc replaced by the retained field proves the
second identity in Eq.~\eqref{eq:pair-jets-zero}.

Now place the field at an internal vertex $r$ of $P_i$.  The raw
$P_i$-field coefficient contains the complete symmetrization of
$S_r^+,S_r^-,S_r^z$.  Its normalized trace is zero because the six traces in
Eq.~\eqref{eq:spin-third-moment} cancel in three opposite pairs.  The field
singleton has trace $\tau_r(S_r^z)=0$, so there is no logarithmic product to
restore the coefficient.  If $j\ne i$, the $P_j$ coefficient factors from
the one-site field trace and vanishes for the same reason.  This proves
Eq.~\eqref{eq:internal-field-pair-zero} for every marked arc.
\end{proof}

\subsection{Composition with the modular functional}

We need two elementary derivatives of the state functional at the maximally
mixed state.  We complexify its real-analytic derivatives when they act on
the formal non-Hermitian oriented arc jets.

\begin{lemma}[First two derivatives]
\label{lem:triangle-functional-low-derivatives}
The functional in Eq.~\eqref{eq:triangle-functional} satisfies

\begin{equation}
 D\mathcal F(0)=0,
 \qquad
 D^2\mathcal F(0)=0.
 \label{eq:F-first-two-zero}
\end{equation}
\end{lemma}

\begin{proof}
For a perturbation $R$, each reduced logarithm is a scalar plus the linear
term $-\E_X(R)$ and terms of higher degree, where $\E_X$ is the normalized
conditional expectation onto $X$.  Their commutator is
therefore at least quadratic, which gives $D\mathcal F(0)=0$.  At quadratic
order, the state in the expectation value may be replaced by the maximally
mixed state.  For two perturbations $R,S$, polarization leaves a sum of
normalized traces of commutators such as
$\ii\tau_U([\E_{AB}(R),\E_{BC}(S)])$.  Every such trace is
zero.  The linear correction to the Gibbs weight first multiplies a
quadratic commutator at cubic order.  Hence $D^2\mathcal F(0)=0$.
\end{proof}

We can now finish the proof of Theorem~\ref{thm:unicyclic-onset}.

\begin{proof}
The grading argument already proves
Eq.~\eqref{eq:unicyclic-low-orders-zero} and reduces the coefficient at order
$n+1$ to one cycle orientation and one field.  By pruning, it remains to
evaluate the isolated cycle.

First suppose that the field lies at one of the retained vertices.  Regard
the individual edge variables and the field variable as the variables of
the composition $\mathcal F(Q)$.  The multiaffine chain rule for the
primitive coefficient is a sum over labeled set partitions, with no
additional factorial.  Lemma
\ref{lem:proper-subpath} says that any nonzero block containing arc edges
must contain one or more complete arcs.  We may therefore organize the
partitions by the four ingredients

\begin{equation}
 \{P_1,P_2,P_3,h_r\}.
 \label{eq:four-ingredients}
\end{equation}

Four singleton blocks contribute the polarized fourth derivative of
$\mathcal F$.  A pair and two singletons would contribute a third derivative,
but every possible pair jet of $Q$ vanishes by
Lemma~\ref{lem:pairwise-shuffle}.  A triple and a singleton, or two pairs,
can contribute only through $D^2\mathcal F(0)$.  One block can contribute
only through $D\mathcal F(0)$.  Lemma
\ref{lem:triangle-functional-low-derivatives} removes the last two cases.
Only the four singleton arc and field jets survive in this primitive
multiaffine coefficient.

If the field lies in the interior of an arc, its singleton jet vanishes under
the partial trace.  A block containing the field and a proper part of that
arc vanishes by Lemma~\ref{lem:proper-subpath}.  The block containing the
field and any complete arc has zero pair jet by
Eq.~\eqref{eq:internal-field-pair-zero}.  Every larger block meets either
$D^2\mathcal F(0)$ or $D\mathcal F(0)$.  Thus the coefficient of
$h_rW_\gamma$ in $m_{n+1}$ vanishes for every internal-cycle vertex $r$.
The pruning argument gives the same conclusion for every off-cycle vertex.
Of the retained fields, Proposition
\ref{prop:arbitrary-spin-triangle} retains only $h_B$.

For an arc $P:u\to v$ of length $\ell$, its singleton jet can be written

\begin{equation}
 [P]Q
 =\frac{\beta}{2}t_{uv}^{\mathrm{eff}}S_u^+S_v^-,
 \qquad
 t_{uv}^{\mathrm{eff}}
 =(-\beta)^{\ell-1}
 \left(\prod_{x\in P^\circ}\kappa_x\right)T_P.
 \label{eq:primitive-effective-hopping}
\end{equation}

Denote the reverse-oriented arc sector by $\overline P$.  On the physical
parameter slice it supplies the adjoint singleton, and
\begin{equation}
 [P]Q+[\overline P]Q
 =\frac{\beta}{2}\left(
 t_{uv}^{\mathrm{eff}}S_u^+S_v^-
 +\overline{t_{uv}^{\mathrm{eff}}}S_u^-S_v^+
 \right).
 \label{eq:primitive-effective-Hermitian-link}
\end{equation}
These are identities only for primitive singleton jets of the dimensionless
operator $Q$.  No effective-Hamiltonian replacement is being made.  The
field singleton is $-\beta h_rS_r^z$.  Polynomial complexification allows
the polarized triangular coefficient in
Proposition~\ref{prop:arbitrary-spin-triangle} to be applied orientation by
orientation.  The conjugate pair then recombines into
$\operatorname{Im}W_\gamma$.  The powers of $\beta$ are
\begin{equation}
 \beta^4\prod_{i=1}^3\beta^{\ell_i-1}=\beta^{n+1},
 \label{eq:three-arc-beta-power}
\end{equation}
and the effective link factors give the sign

\begin{equation}
 (-1)^{(\ell_1-1)+(\ell_2-1)+(\ell_3-1)}=(-1)^{n-3}.
 \label{eq:three-arc-sign}
\end{equation}

The three arc products satisfy
$T_{P_1}T_{P_2}T_{P_3}=W_\gamma$.  Their internal spin factors, together with
$\kappa_A\kappa_B\kappa_C$ from the triangular seed, give
$\prod_{x\in V(\gamma)}\kappa_x$.  Finally,

\begin{equation}
 -2(-1)^{n-3}=2(-1)^n.
 \label{eq:unicyclic-final-sign}
\end{equation}

This proves the coefficient in Eq.~\eqref{eq:unicyclic-onset}.  Analyticity
supplies the stated remainder for fixed Hamiltonian parameters.  The proof
was polynomial after formal complexification, so it also covers vanishing
hoppings and hence $W_\gamma=0$.
\end{proof}

\subsection{Orientation and nonconsecutive marked sites}

No step above assumes that the marked vertices are consecutive.  The three
arc lengths enter the displayed leading coefficient only through their sum,
the cycle product, and the local $\kappa_x$ factors.  Thus the coefficient in
Eq.~\eqref{eq:unicyclic-onset} is independent of the distribution of the
three marked-site separations.

If an external orientation of $\gamma$ is fixed instead of the convention in
Eq.~\eqref{eq:marked-cycle-order}, define

\begin{equation}
 \varepsilon_\gamma(A,B,C)=
 \begin{cases}
 +1 & \text{if $A,B,C$ occur in that cyclic order},\\
 -1 & \text{if $A,C,B$ occur in that cyclic order}.
 \end{cases}
 \label{eq:cyclic-order-sign}
\end{equation}

The right-hand side of Eq.~\eqref{eq:unicyclic-onset} then acquires the factor
$\varepsilon_\gamma(A,B,C)$.  Equivalently, one may always orient the cycle
through $A,B,C$ as in Eq.~\eqref{eq:marked-cycle-order}.  Exchanging $A$ and
$C$ reverses the result, in agreement with the antisymmetry of the modular
commutator.

When $s_x=1/2$ at every cycle vertex, $\kappa_x=1/4$ and

\begin{equation}
 2(-1)^n\prod_{x\in V(\gamma)}\kappa_x
 =\frac{(-1)^n}{2^{2n-1}}.
 \label{eq:qubit-unicyclic-specialization}
\end{equation}

This is the qubit cycle-length law.  The sharper even-$n$ remainder available
for an isolated qubit XY ring is a separate restricted statement.  It is not
part of Theorem~\ref{thm:unicyclic-onset}.
